\subsection{Exponential test spaces and the Fourier boundary}
\label{acfd:global-domain}

The positive prime-power measure has a global distributional realization
on an exponential test space.  This realization specifies the damping
needed before a Schwartz Fourier transform can be taken.
For Fourier transforms, complexify the test-function spaces and use complex-linear
distribution pairings.  The finite arithmetic forms remain real.

For $a\in\mathbb R$, let
\[
 \mathcal S_a(\mathbb R)=
 \{f\in C^\infty(\mathbb R;\mathbb C):e^{at}f(t)\in\mathcal S(\mathbb R)\}.
\]
Give this space the Fr\'echet topology transported by the isomorphism
$M_a:\mathcal S_a\to\mathcal S$, $M_af(t)=e^{at}f(t)$.
Thus its seminorms are
$\sup_t(1+|t|)^r|\partial_t^j(e^{at}f(t))|$, for integers $r,j\geq0$.
Let $\mathcal S_a'$ denote its continuous complex-linear dual.
Compact smooth functions form a dense subspace.

\begin{theorem}[The damped prime-power distribution]\label{acfd:exponential-domain}
The positive prime-power distribution $\mu_\sigma$ has a continuous
extension to $\mathcal S_a$ if and only if $\sigma+a\geq1$.
When this condition holds, the extension is unique and
\begin{equation}\label{acfd:typed-damping}
 \mu_\sigma=M_a'\mu_{\sigma+a},\qquad
 \mathcal J_a\mu_\sigma=\mu_{\sigma+a},\qquad
 (\mathcal J_a T)(g)=T(e^{-at}g).
\end{equation}
Here $M_a':\mathcal S'\to\mathcal S_a'$ is the transpose and
$\mathcal J_a:\mathcal S_a'\to\mathcal S'$ is its inverse.
For $C_*=e\sum_{j\geq1}j^{-2}$, one has
\[
 |\mu_\sigma(f)|\leq
 C_*\sup_t(1+|t|)^3|e^{at}f(t)|.
\]
\end{theorem}

\begin{proof}
On compact smooth functions, the identity
\[
 \mu_\sigma(e^{-at}g)=\mu_{\sigma+a}(g)
\]
holds term by term in the locally finite atomic sums.
If $\sigma+a\geq1$, Theorem~\ref{acfd:tempered} proves the estimate and
therefore the continuous extension through $M_a$.
Density gives uniqueness.
Conversely, a continuous extension on $\mathcal S_a$ would give a tempered
extension of $\mu_{\sigma+a}$ by composition with $M_a^{-1}$.
The same theorem excludes this when $\sigma+a<1$.
\end{proof}

For $g\in\mathcal S(\mathbb R)$, fix the Fourier convention
\[
 \widehat g(\xi)=\int_{\mathbb R}e^{-it\xi}g(t)\,dt,
 \qquad
 \widehat T(g)=T(\widehat g)\quad(T\in\mathcal S').
\]
The inverse Fourier transform on functions has factor $(2\pi)^{-1}$
and exponential $e^{it\xi}$.
Repeated differentiation and integration by parts prove continuity of the
Fourier map on $\mathcal S$.  More explicitly,
\[
 \|\xi^r\partial_\xi^j\widehat g\|_\infty
    \leq\|\partial_t^r(t^jg)\|_{L^1}
       \qquad(r,j\in\mathbb Z_{\geq0}).
\]
The product rule bounds the right side by a finite sum of Schwartz
seminorms, after integration against $(1+|t|)^{-2}$.

\begin{theorem}[A distributional boundary of the logarithmic derivative]
\label{acfd:fourier-boundary}
For real $s>1$, define $\zeta(z)=\sum_{n\geq1}n^{-z}$ on
$\operatorname{Re}z>1$.  Then $\widehat\mu_s$ is represented by the smooth function
\begin{equation}\label{acfd:dirichlet-fourier}
 G_s(\xi)=\sum_{p,m\geq1}(\log p)p^{-m(s+i\xi)}
              =-\frac{\zeta'(s+i\xi)}{\zeta(s+i\xi)}.
\end{equation}
Every derivative series converges uniformly in $\xi\in\mathbb R$.
As $\epsilon\downarrow0$,
\begin{equation}\label{acfd:boundary-limit}
 G_{1+\epsilon}\longrightarrow\widehat\mu_1
                 \quad\text{in the strong topology of }\mathcal S',
 \qquad
 |(\widehat\mu_{1+\epsilon}-\widehat\mu_1)(g)|
       \leq C_*\epsilon\sup_t(1+|t|)^4|\widehat g(t)|.
\end{equation}
In particular, for the critical weight $\sigma=1/2$, there is the typed sequence
\[
 \mu_{1/2}\in\mathcal S_{1/2}'
 \xrightarrow{\ \mathcal J_{1/2}\ }
 \mu_1\in\mathcal S'
 \xrightarrow{\ \mathcal F\ }
 \widehat\mu_1\in\mathcal S'.
\]
The first distribution belongs to the exponential dual stated here.
It has no extension to the unweighted Schwartz space.
\end{theorem}

\begin{proof}
Unique factorization assigns at most one prime base to each prime power.
For every nonnegative integer $j$,
\[
 \sum_{p,m\geq1}(\log p)(m\log p)^j p^{-ms}
       \leq\sum_{n\geq2}\frac{(\log n)^{j+1}}{n^s}<\infty.
\]
The last convergence follows by an integral comparison.
This proves the derivative convergence and justifies Fubini in the Fourier pairing.

For $\operatorname{Re}z>1$, absolute convergence and unique factorization give
\[
 \zeta(z)=\prod_p(1-p^{-z})^{-1}
     =\exp\left(\sum_p\sum_{m\geq1}\frac{p^{-mz}}m\right).
\]
The last sum and its derivative converge locally uniformly in that half-plane.
The exponential is nonzero.  Differentiation proves
\eqref{acfd:dirichlet-fourier}.

For $\epsilon>0$, the positive measures satisfy
$d\mu_{1+\epsilon}(t)=e^{-\epsilon t}d\mu_1(t)$ and have support in $t>0$.
The inequality $1-e^{-\epsilon t}\leq\epsilon t$ gives
\[
 |(\mu_{1+\epsilon}-\mu_1)(u)|
       \leq\epsilon\mu_1(t|u(t)|)
       \leq C_*\epsilon\sup_t(1+|t|)^4|u(t)|.
\]
Substitute $u=\widehat g$ to obtain \eqref{acfd:boundary-limit}.
Every bounded subset of $\mathcal S$ has bounded Fourier seminorms.
The same estimate therefore proves strong convergence in its continuous dual.
The typed sequence follows from Theorem~\ref{acfd:exponential-domain}.
\end{proof}

The pole at frequency zero gives a direct normalization check for this boundary.
For $\operatorname{Re}z>1$, write
\[
 \zeta(z)=\frac1{z-1}+H(z),\qquad
 H(z)=\sum_{n\geq1}
        \left(n^{-z}-\int_n^{n+1}t^{-z}\,dt\right).
\]
For $z$ in a compact subset of $\operatorname{Re}z>0$, each summand is
bounded by a constant times $n^{-\operatorname{Re}z-1}$.
Indeed, integrate the derivative $-zt^{-z-1}$ between $n$ and $t$.
Thus $H$ is holomorphic on that half-plane by locally uniform convergence.
Near $z=1$, the function $1+(z-1)H(z)$ is nonzero.  Consequently
\[
 -\frac{\zeta'(z)}{\zeta(z)}
       =\frac1{z-1}+A(z)
\]
for a holomorphic function $A$ in a neighborhood of one.

\begin{proposition}[The zero-frequency model]\label{acfd:pole-model}
On some interval $(-r,r)$, the distribution $\widehat\mu_1$ is
\begin{equation}\label{acfd:pole-sign}
 \widehat\mu_1
      =\pi\delta_0-i\,\operatorname{pv}\frac1\xi+A(1+i\xi).
\end{equation}
The principal value is defined on compact smooth tests by symmetric
removal of $(-b,b)$ followed by $b\downarrow0$.
\end{proposition}

\begin{proof}
Choose $r$ so that $1+(z-1)H(z)$ is nonzero for $|z-1|<2r$.
For a test $g\in C_c^\infty((-r,r))$, the holomorphic term converges
uniformly on its support as $z=1+\epsilon+i\xi$ approaches the line.
The pole is
\[
 \frac1{\epsilon+i\xi}
   =\frac{\epsilon}{\epsilon^2+\xi^2}
       -i\frac{\xi}{\epsilon^2+\xi^2}.
\]
The first kernel has integral $\pi$, concentrates at zero, and converges
to $\pi\delta_0$.
For the second, subtract $g(0)$ times an even compact cutoff equal to one
near zero.  Its pairing with the odd kernel is zero.
The remaining test divided by $\xi$ is bounded near zero, so dominated
convergence gives $\operatorname{pv}(1/\xi)$ with the displayed sign.
Combine this calculation with \eqref{acfd:boundary-limit}.
\end{proof}

This construction gives a global distribution and a local pole calculation
for the arithmetic prime-power source.
Comparison with the simplex operator $K_k$ requires an identity for
bilinear forms on specified carriers.
The following construction gives such a pairing through finite incidence.
A complete explicit formula also needs its other terms and a common domain
on which their sum is defined.
